% GENERATED FILE. Edit canonical lessons, then run npm run generate:books.
% canonical-source-hash: fnv1a64:a1aa6c19a32832e3
% canonical-lessons: HI-C456-naraz, HI-C456-krodh, HI-C456-cirhna, HI-C456-bipharna, HI-C456-khijhna

\chapter{Losing One's Temper}
\label{ch:hi-losing-one-s-temper}
\hlchaptermodality{456}

\begin{chapteropening}
\textbf{By the end of this chapter:} \emph{I can say that someone got cross and why.}

Complete the last lesson of chapter 456: I can say that someone got cross and why.
\end{chapteropening}

\section[\texorpdfstring{\hi{नाराज़} (nārāz)}{नाराज़ (nārāz)}]{\hi{नाराज़} (nārāz) — cross; displeased}

\label{lesson:HI-C456-naraz}



\begin{quote}
Before the new one: say the Hindi for to fret, then the Hindi for to hurt someone's feelings.
\end{quote}

\subsection*{You'll want to know: \hi{नाराज़}}
\textbf{\hi{नाराज़}} — \emph{nārāz} — \textquotedblleft{}cross; displeased\textquotedblright{}.

\begin{quote}
\hi{पिता} \hi{अरुण} \hi{से} \hi{नाराज़} \hi{थे} \hi{क्योंकि} \hi{वह} \hi{देर} \hi{से} \hi{आया।}
\end{quote}

\begin{quote}
\emph{pitā Aruṇ se nārāz the kyoṁki vah der se āyā.}
\end{quote}

\begin{quote}
Arun's father was cross with him because he came late.
\end{quote}

\begin{cousinweb}
\textbf{\hi{नाराज़}} (\emph{nārāz}) is usually traced to Persian \emph{nārāz}, \textquotedblleft{}displeased\textquotedblright{}, built on \emph{nā-}, \textquotedblleft{}not\textquotedblright{}, and a word for \textquotedblleft{}content\textquotedblright{}.
\end{cousinweb}

\begin{grammarlens}[title={What you've built}]
One more word for describing what you lived and felt.
\end{grammarlens}

\subsection*{Your turn}
\begin{itemize}
\raggedright
  \item \emph{Say it:} \emph{nārāz}
  \item \emph{Say it:} \emph{nārāz}, once more
  \item \emph{Say it:} \emph{pitā Aruṇ se nārāz the kyoṁki vah der se āyā.} — Arun's father was cross with him because he came late.
  \item \emph{Recall:} say the Hindi for to fret, then the Hindi for to hurt someone's feelings, then say \emph{nārāz} again
\end{itemize}

\begin{tcolorbox}[breakable,colback=teal!4,colframe=teal!35!black,title={Before you move on}]
What does \hi{नाराज़} mean? (\textquotedblleft{}Cross; displeased\textquotedblright{}.) Say it once more.
\end{tcolorbox}
\section[\texorpdfstring{\hi{क्रोध} (krodh)}{क्रोध (krodh)}]{\hi{क्रोध} (krodh) — wrath; rage}

\label{lesson:HI-C456-krodh}



\begin{quote}
Before the new one: say the Hindi for to hurt someone's feelings, then the Hindi for cross.
\end{quote}

\subsection*{You'll want to know: \hi{क्रोध}}
\textbf{\hi{क्रोध}} — \emph{krodh} — \textquotedblleft{}wrath; rage\textquotedblright{}.

\begin{quote}
\hi{क्रोध} \hi{में} \hi{आदमी} \hi{सही} \hi{और} \hi{ग़लत} \hi{भूल} \hi{जाता} \hi{है।}
\end{quote}

\begin{quote}
\emph{krodh meṁ ādmī sahī aur ġalat bhūl jātā hai.}
\end{quote}

\begin{quote}
In a rage a person forgets right from wrong.
\end{quote}

Masculine.

\begin{cousinweb}
\textbf{\hi{क्रोध}} (\emph{krodh}) is a tatsama, borrowed from Sanskrit \emph{krodha}, \textquotedblleft{}anger, wrath, passion\textquotedblright{} (Monier-Williams). Everyday speech says \emph{ġussā}; \emph{krodh} is the word of books and sermons.
\end{cousinweb}

\begin{grammarlens}[title={What you've built}]
One more word for describing what you lived and felt.
\end{grammarlens}

\subsection*{Your turn}
\begin{itemize}
\raggedright
  \item \emph{Say it:} \emph{krodh}
  \item \emph{Say it:} \emph{krodh}, once more
  \item \emph{Say it:} \emph{krodh meṁ ādmī sahī aur ġalat bhūl jātā hai.} — In a rage a person forgets right from wrong.
  \item \emph{Recall:} say the Hindi for to hurt someone's feelings, then the Hindi for cross, then say \emph{krodh} again
\end{itemize}

\begin{tcolorbox}[breakable,colback=teal!4,colframe=teal!35!black,title={Before you move on}]
What does \hi{क्रोध} mean? (\textquotedblleft{}Wrath; rage\textquotedblright{}.) Say it once more.
\end{tcolorbox}
\section[\texorpdfstring{\hi{चिढ़ना} (ciṛhnā)}{चिढ़ना (ciṛhnā)}]{\hi{चिढ़ना} (ciṛhnā) — to get irritated}

\label{lesson:HI-C456-cirhna}



\begin{quote}
Before the new one: say the Hindi for cross, then the Hindi for wrath.
\end{quote}

\subsection*{You'll want to know: \hi{चिढ़ना}}
\textbf{\hi{चिढ़ना}} — \emph{ciṛhnā} — \textquotedblleft{}to get irritated\textquotedblright{}.

\begin{quote}
\hi{मीरा} \hi{का} \hi{भाई} \hi{छोटी} \hi{बात} \hi{पर} \hi{चिढ़} \hi{जाता} \hi{है।}
\end{quote}

\begin{quote}
\emph{Mīrā kā bhāī choṭī bāt par ciṛh jātā hai.}
\end{quote}

\begin{quote}
Mira's brother gets irritated over little things.
\end{quote}

Past (he …): \hi{वह} \hi{चिढ़} \hi{गया}, \emph{vah ciṛh gayā}.

\begin{cousinweb}
\textbf{\hi{चिढ़ना}} (\emph{ciṛhnā}) is a tadbhava: CDIAL lists Hindi \emph{ciṛhnā}, \textquotedblleft{}to be provoked\textquotedblright{}, under \emph{*ciḍ}, a form it reconstructs (CDIAL 4794).
\end{cousinweb}

\begin{grammarlens}[title={What you've built}]
One more word for describing what you lived and felt.
\end{grammarlens}

\subsection*{Your turn}
\begin{itemize}
\raggedright
  \item \emph{Say it:} \emph{ciṛhnā}
  \item \emph{Say it:} \emph{ciṛhnā}, once more
  \item \emph{Say it:} \emph{Mīrā kā bhāī choṭī bāt par ciṛh jātā hai.} — Mira's brother gets irritated over little things.
  \item \emph{Recall:} say the Hindi for cross, then the Hindi for wrath, then say \emph{ciṛhnā} again
\end{itemize}

\begin{tcolorbox}[breakable,colback=teal!4,colframe=teal!35!black,title={Before you move on}]
What does \hi{चिढ़ना} mean? (\textquotedblleft{}To get irritated\textquotedblright{}.) Say it once more.
\end{tcolorbox}
\section[\texorpdfstring{\hi{बिफरना} (bipharnā)}{बिफरना (bipharnā)}]{\hi{बिफरना} (bipharnā) — to fly into a rage}

\label{lesson:HI-C456-bipharna}



\begin{quote}
Before the new one: say the Hindi for wrath, then the Hindi for to get irritated.
\end{quote}

\subsection*{You'll want to know: \hi{बिफरना}}
\textbf{\hi{बिफरना}} — \emph{bipharnā} — \textquotedblleft{}to fly into a rage\textquotedblright{}.

\begin{quote}
\hi{कम} \hi{पैसे} \hi{देखकर} \hi{दुकानदार} \hi{बिफर} \hi{गया।}
\end{quote}

\begin{quote}
\emph{kam paise dekhkar dukāndār biphar gayā.}
\end{quote}

\begin{quote}
Seeing how little money he had been given, the shopkeeper flew into a rage.
\end{quote}

Past (he …): \hi{वह} \hi{बिफर} \hi{गया}, \emph{vah biphar gayā}.

\begin{cousinweb}
\textbf{\hi{बिफरना}} (\emph{bipharnā}) is a tadbhava: CDIAL lists Hindi \emph{bipharnā}, \textquotedblleft{}to become enraged\textquotedblright{}, under \emph{*vispharati}, \textquotedblleft{}quivers, trembles\textquotedblright{}, a form it reconstructs (CDIAL 12014): someone shaking with anger. It is an outburst, much stronger than \emph{ciṛhnā}.
\end{cousinweb}

\begin{grammarlens}[title={What you've built}]
One more word for describing what you lived and felt.
\end{grammarlens}

\subsection*{Your turn}
\begin{itemize}
\raggedright
  \item \emph{Say it:} \emph{bipharnā}
  \item \emph{Say it:} \emph{bipharnā}, once more
  \item \emph{Say it:} \emph{kam paise dekhkar dukāndār biphar gayā.} — Seeing how little money he had been given, the shopkeeper flew into a rage.
  \item \emph{Recall:} say the Hindi for wrath, then the Hindi for to get irritated, then say \emph{bipharnā} again
\end{itemize}

\begin{tcolorbox}[breakable,colback=teal!4,colframe=teal!35!black,title={Before you move on}]
What does \hi{बिफरना} mean? (\textquotedblleft{}To fly into a rage\textquotedblright{}.) Say it once more.
\end{tcolorbox}
\section[\texorpdfstring{\hi{खीझना} (khījhnā)}{खीझना (khījhnā)}]{\hi{खीझना} (khījhnā) — to get fed up; to be vexed}

\label{lesson:HI-C456-khijhna}



\begin{quote}
Before the new one: say the Hindi for to get irritated, then the Hindi for to fly into a rage.
\end{quote}

\subsection*{You'll want to know: \hi{खीझना}}
\textbf{\hi{खीझना}} — \emph{khījhnā} — \textquotedblleft{}to get fed up; to be vexed\textquotedblright{}.

\begin{quote}
\hi{रात} \hi{को} \hi{फ़ोन} \hi{तीन} \hi{बार} \hi{बजा} \hi{और} \hi{मीरा} \hi{खीझ} \hi{गई।}
\end{quote}

\begin{quote}
\emph{rāt ko fon tīn bār bajā aur Mīrā khījh gaī.}
\end{quote}

\begin{quote}
The phone rang three times in the night and Mira got fed up.
\end{quote}

Past (she …): \hi{वह} \hi{खीझ} \hi{गई}, \emph{vah khījh gaī}.

\begin{cousinweb}
\textbf{\hi{खीझना}} (\emph{khījhnā}) is inherited: CDIAL lists Hindi \emph{khījnā}, \textquotedblleft{}to be irritated\textquotedblright{}, and the causative \emph{khijhānā}, \textquotedblleft{}to vex\textquotedblright{}, under Sanskrit \emph{kṣīyátē}, \textquotedblleft{}is destroyed, wanes\textquotedblright{} (CDIAL 3695); \emph{khījhnā} is the usual spelling of that plain verb.
\end{cousinweb}

\begin{grammarlens}[title={What you've built}]
One more word for describing what you lived and felt.
\end{grammarlens}

\subsection*{Your turn}
\begin{itemize}
\raggedright
  \item \emph{Say it:} \emph{khījhnā}
  \item \emph{Say it:} \emph{khījhnā}, once more
  \item \emph{Say it:} \emph{rāt ko fon tīn bār bajā aur Mīrā khījh gaī.} — The phone rang three times in the night and Mira got fed up.
  \item \emph{Recall:} say the Hindi for to get irritated, then the Hindi for to fly into a rage, then say \emph{khījhnā} again
\end{itemize}

\begin{tcolorbox}[breakable,colback=teal!4,colframe=teal!35!black,title={Before you move on}]
What does \hi{खीझना} mean? (\textquotedblleft{}To get fed up; to be vexed\textquotedblright{}.) Say it once more.
\end{tcolorbox}
